\documentclass[11pt]{article}
\usepackage[T1]{fontenc}
\usepackage{lmodern}
\usepackage{microtype}
\usepackage[margin=0.68in,headheight=15pt]{geometry}
\usepackage{amsmath,amssymb,mathtools}
\usepackage{xcolor}
\usepackage[most]{tcolorbox}
\usepackage{enumitem}
\usepackage{booktabs,array}
\usepackage{fancyhdr}
\usepackage{hyperref}

\definecolor{indigo}{HTML}{25355D}
\definecolor{violet}{HTML}{6857B2}
\definecolor{mint}{HTML}{2B9B82}
\definecolor{amber}{HTML}{D99A28}
\definecolor{rose}{HTML}{CF5D75}
\definecolor{paper}{HTML}{FAF8F3}
\definecolor{ink}{HTML}{222834}
\definecolor{muted}{HTML}{687080}

\hypersetup{colorlinks=true,linkcolor=violet,urlcolor=violet}
\pagestyle{fancy}
\fancyhf{}
\renewcommand{\headrulewidth}{0pt}
\fancyfoot[L]{\footnotesize\color{muted}Chance, Data, and Decisions}
\fancyfoot[R]{\footnotesize\color{muted}Quiz \thequiznum\ of 10}

\setlength{\parindent}{0pt}
\setlength{\parskip}{3pt}
\setlist[enumerate]{leftmargin=*,itemsep=3pt,topsep=4pt}
\newcounter{quiznum}
\newcommand{\E}{\mathbb E}
\newcommand{\Var}{\operatorname{Var}}
\newcommand{\Cov}{\operatorname{Cov}}
\newcommand{\Prob}{\mathbb P}
\newcommand{\edition}{\ifanswers Instructor Key\else Student Edition\fi}
\newcommand{\points}[1]{\hfill{\small\color{rose}\textbf{[#1 points]}}}

\newcommand{\quizheading}[1]{%
  \clearpage\stepcounter{quiznum}%
  \begin{tcolorbox}[
    enhanced,colback=paper,colframe=indigo,boxrule=1.1pt,arc=2mm,
    left=4mm,right=4mm,top=2.5mm,bottom=2.5mm,
    borderline west={4pt}{0pt}{violet}]
    {\Large\bfseries\color{indigo}Weekly Quiz \thequiznum}\hfill
    {\bfseries\color{violet}\edition}\\[-1pt]
    {\large\color{ink}#1}\\[2pt]
    {\small\color{muted}20 minutes \quad $\bullet$ \quad 10 points \quad $\bullet$ \quad Show the essential steps.}
    \ifanswers\else
      \\[5pt]\textbf{Name:}\ \rule{0.48\linewidth}{0.4pt}\hfill
      \textbf{Date:}\ \rule{0.20\linewidth}{0.4pt}
    \fi
  \end{tcolorbox}
}

\NewDocumentEnvironment{solution}{O{1.05in}+b}{%
  \ifanswers
    \begin{tcolorbox}[enhanced,breakable,colback=mint!4,colframe=mint!45,
      boxrule=0.55pt,arc=1mm,left=2.5mm,right=2.5mm,top=1.5mm,bottom=1.5mm,
      title={\small\bfseries Solution},fonttitle=\small\bfseries,coltitle=indigo]
      \small #2
    \end{tcolorbox}
  \else
    \par\vspace{#1}
  \fi
}{}

\begin{document}

\quizheading{Sets, diagrams, and counting}
\begin{enumerate}
  \item Express ``$A$ or $B$, but not both'' using unions, intersections, and
  complements. \points{3}
  \begin{solution}[1.05in]
    The event is the symmetric difference
    \[
       (A\cap B^c)\cup(A^c\cap B),
    \]
    equivalently $(A\cup B)\setminus(A\cap B)$.
  \end{solution}

  \item A fair coin is tossed three times. Find the probability of exactly two
  heads. \points{3}
  \begin{solution}[1.2in]
    There are $2^3=8$ equally likely strings. Choose the two positions occupied by
    heads in $\binom32=3$ ways, so the probability is $3/8$.
  \end{solution}

  \item In a group of $70$ students, $38$ study French, $27$ study Spanish, and
  $12$ study both. How many study neither language? \points{4}
  \begin{solution}[1.45in]
    Inclusion--exclusion gives
    \[
      |F\cup S|=|F|+|S|-|F\cap S|=38+27-12=53.
    \]
    Therefore $70-53=17$ students study neither.
  \end{solution}
\end{enumerate}

\quizheading{Probability spaces and events}
\begin{enumerate}
  \item Two fair dice are rolled. Find the probability that the maximum is $4$.
  \points{4}
  \begin{solution}[1.4in]
    There are $4^2=16$ outcomes with both dice at most $4$, and $3^2=9$ with both
    at most $3$. Thus $16-9=7$ of the $36$ outcomes have maximum $4$, giving
    probability $7/36$.
  \end{solution}

  \item Find the probability of at least one six in four independent rolls of a
  fair die. \points{3}
  \begin{solution}[1.15in]
    The complement is no six in all four rolls, which has probability $(5/6)^4$.
    Therefore
    \[
       \Prob(\text{at least one six})=1-\left(\frac56\right)^4=\frac{671}{1296}.
    \]
  \end{solution}

  \item Derive $\Prob(A\cup B)=\Prob(A)+\Prob(B)-\Prob(A\cap B)$ by splitting
  events into disjoint pieces. \points{3}
  \begin{solution}[1.35in]
    The union is the disjoint union $A\cup(B\setminus A)$, so
    $\Prob(A\cup B)=\Prob(A)+\Prob(B\setminus A)$. Also
    $B=(B\setminus A)\cup(A\cap B)$ disjointly, hence
    $\Prob(B\setminus A)=\Prob(B)-\Prob(A\cap B)$. Substitution gives the formula.
  \end{solution}
\end{enumerate}

\quizheading{Conditional probability and Bayes' rule}
\begin{enumerate}
  \item An urn contains $3$ red and $2$ blue balls. Two are drawn without
  replacement. Find the probability that the second ball is red. \points{3}
  \begin{solution}[1.25in]
    By conditioning on the first draw,
    \[
    \Prob(R_2)=\frac35\frac24+\frac25\frac34=\frac{6+6}{20}=\frac35.
    \]
    (Symmetry also says every draw position is red with probability $3/5$.)
  \end{solution}

  \item A condition has prevalence $2\%$. A test has sensitivity $90\%$ and
  specificity $95\%$. Find the probability that a person who tests positive has
  the condition. \points{4}
  \begin{solution}[1.55in]
    For a random person,
    \[
      \Prob(D\cap +)=0.02(0.90)=0.018,
    \]
    while the false-positive probability is $0.98(0.05)=0.049$. Bayes' rule gives
    \[
      \Prob(D\mid +)=\frac{0.018}{0.018+0.049}=\frac{18}{67}\approx0.269.
    \]
  \end{solution}

  \item Suppose $\Prob(A)=0.4$, $\Prob(B)=0.5$, and $\Prob(A\cap B)=0.2$.
  Are $A$ and $B$ independent? Find $\Prob(A\mid B)$. \points{3}
  \begin{solution}[1.2in]
    Since $\Prob(A)\Prob(B)=0.4(0.5)=0.2=\Prob(A\cap B)$, the events are
    independent. Also $\Prob(A\mid B)=0.2/0.5=0.4$.
  \end{solution}
\end{enumerate}

\quizheading{Random variables, PMFs, densities, and CDFs}
\begin{enumerate}
  \item A random variable $X$ takes values $0,1,2$ with probabilities
  $0.2,0.5,0.3$. Write its CDF $F(x)=\Prob(X\le x)$ piecewise. \points{4}
  \begin{solution}[1.55in]
    \[
    F(x)=
    \begin{cases}
      0,&x<0,\\
      0.2,&0\le x<1,\\
      0.7,&1\le x<2,\\
      1,&x\ge2.
    \end{cases}
    \]
  \end{solution}

  \item Let $f(x)=cx$ for $0\le x\le2$ and $f(x)=0$ otherwise. Find $c$ and
  $\Prob(X>1)$. \points{4}
  \begin{solution}[1.45in]
    Normalization gives $1=\int_0^2cx\,dx=2c$, so $c=1/2$. Then
    \[
      \Prob(X>1)=\int_1^2\frac{x}{2}\,dx
      =\left.\frac{x^2}{4}\right|_1^2=\frac34.
    \]
  \end{solution}

  \item If $X$ has a continuous distribution, what is $\Prob(X=1)$? Why?
  \points{2}
  \begin{solution}[0.9in]
    It is $0$: a single point has zero width, so its probability is
    $\int_1^1 f(x)\,dx=0$.
  \end{solution}
\end{enumerate}

\quizheading{Expectation, variance, and moments}
\begin{enumerate}
  \item For $X$ with probabilities $0.2,0.5,0.3$ at $0,1,2$, compute
  $\E[X]$ and $\Var(X)$. \points{4}
  \begin{solution}[1.45in]
    \[
      \E[X]=0(0.2)+1(0.5)+2(0.3)=1.1,
    \]
    and $\E[X^2]=0+1(0.5)+4(0.3)=1.7$. Therefore
    $\Var(X)=\E[X^2]-\E[X]^2=1.7-1.21=0.49$.
  \end{solution}

  \item If $X\sim\operatorname{Bin}(20,0.3)$, find its mean and variance. Explain
  briefly using indicators. \points{3}
  \begin{solution}[1.3in]
    Write $X=I_1+\cdots+I_{20}$, where each $I_i$ is Bernoulli$(0.3)$. Then
    \[
      \E[X]=20(0.3)=6,
      \qquad \Var(X)=20(0.3)(0.7)=4.2,
    \]
    using independence for the variance sum.
  \end{solution}

  \item Given $\E[X]=4$ and $\E[Y]=-1$, find $\E[3X-2Y+5]$. Is independence
  needed? \points{3}
  \begin{solution}[1.1in]
    By linearity,
    $\E[3X-2Y+5]=3(4)-2(-1)+5=19$. Independence is not needed for linearity of
    expectation.
  \end{solution}
\end{enumerate}

\quizheading{Standard probability distributions}
\begin{enumerate}
  \item If $X\sim\operatorname{Bin}(5,0.4)$, find $\Prob(X=2)$. \points{3}
  \begin{solution}[1.15in]
    \[
      \Prob(X=2)=\binom52(0.4)^2(0.6)^3
      =10(0.16)(0.216)=0.3456.
    \]
  \end{solution}

  \item A waiting time is exponential with mean $4$ minutes. Find the probability
  that the wait exceeds $6$ minutes. \points{3}
  \begin{solution}[1.2in]
    The rate is $\lambda=1/4$ per minute. Therefore
    \[
      \Prob(X>6)=e^{-\lambda6}=e^{-3/2}\approx0.223.
    \]
  \end{solution}

  \item Suppose $X\sim N(100,15^2)$. Express $\Prob(X\le130)$ using the standard
  normal CDF $\Phi$, and evaluate using $\Phi(2)\approx0.9772$. \points{4}
  \begin{solution}[1.4in]
    Standardizing,
    \[
      \Prob(X\le130)=\Prob\left(Z\le\frac{130-100}{15}\right)
      =\Prob(Z\le2)=\Phi(2)\approx0.9772.
    \]
  \end{solution}
\end{enumerate}

\quizheading{Joint, marginal, and conditional distributions}
For the next three questions, use the joint probability table
\[
\begin{array}{c|cc|c}
 &Y=0&Y=1&\Prob(X=x)\\ \hline
X=0&0.10&0.20& \\
X=1&0.30&0.40& \\ \hline
\Prob(Y=y)&&&1
\end{array}
\]
\begin{enumerate}
  \item Find both marginal distributions. \points{3}
  \begin{solution}[0.9in]
    Row sums give $\Prob(X=0)=0.30$ and $\Prob(X=1)=0.70$. Column sums give
    $\Prob(Y=0)=0.40$ and $\Prob(Y=1)=0.60$.
  \end{solution}

  \item Find $\Prob(X=1\mid Y=0)$ and decide whether $X$ and $Y$ are independent.
  \points{3}
  \begin{solution}[1.05in]
    $\Prob(X=1\mid Y=0)=0.30/0.40=0.75$. They are not independent: for example,
    $\Prob(X=1,Y=0)=0.30\ne0.70(0.40)=0.28$.
  \end{solution}

  \item Compute $\Cov(X,Y)$. \points{4}
  \begin{solution}[1.15in]
    Since $X,Y$ are indicators, $\E[X]=0.70$, $\E[Y]=0.60$, and
    $\E[XY]=\Prob(X=1,Y=1)=0.40$. Thus
    \[
      \Cov(X,Y)=0.40-(0.70)(0.60)=-0.02.
    \]
  \end{solution}
\end{enumerate}

\quizheading{Sampling distributions, the LLN, and the CLT}
\begin{enumerate}
  \item A population has standard deviation $12$. For samples of size $64$, find
  the standard error of the sample mean. \points{3}
  \begin{solution}[1.0in]
    \[
      \operatorname{SE}(\bar X)=\frac{\sigma}{\sqrt n}
      =\frac{12}{\sqrt{64}}=1.5.
    \]
  \end{solution}

  \item In one sentence each, distinguish the law of large numbers from the central
  limit theorem. \points{3}
  \begin{solution}[1.4in]
    The law of large numbers says the sample mean converges to the population mean
    as the sample grows. The central limit theorem says the standardized sampling
    distribution of the sample mean becomes approximately normal (under its usual
    conditions), describing the shape and scale of its fluctuations.
  \end{solution}

  \item If $100$ independent Bernoulli trials have success probability $0.4$, give
  the approximate mean and standard deviation of the sample proportion. Standardize
  the event $\hat p\le0.50$. \points{4}
  \begin{solution}[1.5in]
    \[
      \E[\hat p]=0.4,\qquad
      \operatorname{SD}(\hat p)=\sqrt{\frac{0.4(0.6)}{100}}\approx0.0490.
    \]
    Thus
    \[
      \Prob(\hat p\le0.50)\approx
      \Phi\left(\frac{0.50-0.40}{0.0490}\right)=\Phi(2.04).
    \]
  \end{solution}
\end{enumerate}

\quizheading{Estimation, confidence intervals, and tests}
\begin{enumerate}
  \item In a random sample, $330$ of $600$ respondents answer yes. Compute the
  usual approximate $95\%$ confidence interval for the population proportion,
  using $z^*=1.96$. \points{4}
  \begin{solution}[1.55in]
    $\hat p=330/600=0.55$ and
    \[
      \operatorname{SE}=\sqrt{\frac{0.55(0.45)}{600}}\approx0.02031.
    \]
    Thus $0.55\pm1.96(0.02031)=0.55\pm0.0398$, or approximately
    $(0.510,0.590)$.
  \end{solution}

  \item Give a correct repeated-sampling interpretation of a $95\%$ confidence
  procedure. \points{3}
  \begin{solution}[1.35in]
    If we repeatedly sampled by the same design and built an interval by this same
    procedure, about $95\%$ of those intervals would contain the fixed population
    parameter. The $95\%$ describes the procedure's long-run coverage, not a random
    probability assigned to the parameter after this interval is observed.
  \end{solution}

  \item A valid test returns $p=0.03$. State an appropriate conclusion at level
  $\alpha=0.05$, and name one conclusion the $p$-value does \emph{not} justify.
  \points{3}
  \begin{solution}[1.3in]
    Since $0.03<0.05$, reject the null hypothesis at the $5\%$ level; the data are
    relatively incompatible with the null model. The result does not say that the
    null hypothesis has probability $0.03$, nor does it by itself establish a large
    or practically important effect.
  \end{solution}
\end{enumerate}

\quizheading{Bayesian inference and prediction}
\begin{enumerate}
  \item A Bernoulli success probability has prior $p\sim\operatorname{Beta}(2,2)$.
  After observing $8$ successes and $2$ failures, find the posterior distribution.
  \points{4}
  \begin{solution}[1.25in]
    Beta--Bernoulli conjugacy adds successes to the first shape and failures to the
    second. Hence
    \[
      p\mid\text{data}\sim\operatorname{Beta}(2+8,2+2)
      =\operatorname{Beta}(10,4).
    \]
  \end{solution}

  \item Find the posterior mean and the posterior predictive probability that the
  next trial is a success. \points{3}
  \begin{solution}[1.15in]
    For $\operatorname{Beta}(10,4)$ the posterior mean is
    $10/(10+4)=5/7$. The posterior predictive success probability is also
    $\E[p\mid\text{data}]=5/7$.
  \end{solution}

  \item Contrast a $95\%$ Bayesian credible interval with a $95\%$ frequentist
  confidence interval in one or two sentences. \points{3}
  \begin{solution}[1.4in]
    Given the prior and model, a $95\%$ credible interval contains $95\%$ of the
    posterior probability for the parameter. A $95\%$ confidence interval comes
    from a procedure that covers the fixed parameter in $95\%$ of repeated samples;
    it does not assign posterior probability without additional assumptions.
  \end{solution}
\end{enumerate}

\end{document}
